% =====================================================================
%  PLAN DE TRABAJO (PT) -- Concurso LE-403/729 (UPC)
%  Borrador generado el 05/11/2026. Docencia: solo asignaturas de la convocatoria.
%  Investigación: solo las tres líneas de investigación de la convocatoria.
%  Completar cada \TODO{}; activar \finaltrue antes de presentar.
%  Nombre del archivo de salida: Monteiro_Carlos_PT.pdf
% =====================================================================
\documentclass[11pt]{article}
\usepackage{lector}
\usepackage{pgfgantt}
%\finaltrue
\setupheader{PT}

\title{\textbf{Plan de Trabajo}}
\author{\Large \CandName}
\date{}

\begin{document}
\maketitle
\thispagestyle{plain}
\vspace{0.5cm}
\begin{center}\postdata\end{center}

\vspace{0.8cm}
\noindent\textbf{Resumen.}
Este plan de trabajo abarca los seis años del contrato de profesor lector (2027--2033),
con una autoevaluación al final del tercer año que coincide con la evaluación prevista en
el artículo 78 de la LOSU. En investigación, consolidaré una línea sobre \emph{turbulencia
multifísica en fluidos compresibles y supercríticos}, analizada mediante \emph{métodos
basados en datos y de causalidad} para la predicción de eventos extremos, y calculada con
\emph{solvers distribuidos y acelerados en GPU}. Estas son las tres líneas de investigación
de la convocatoria. En docencia, contribuiré a las asignaturas termofluidodinámicas de la
EEBE: Termodinámica y Transferencia de Calor (teoría), Mecánica de Fluidos (laboratorio),
Ingeniería de Fluidos (teoría) y la asignatura de máster Mecánica de Fluidos Computacional.
Introduciré metodologías computacionales y de aprendizaje activo, y vincularé los trabajos
de fin de estudios al plan de investigación.

\newpage
\tableofcontents
\newpage

% =====================================================================
\section{Proyecto de investigación}

\subsection{Antecedentes}
Mi tesis doctoral en la UPC (2023--2026, Sobresaliente \emph{cum laude}) estudió la
\emph{turbulencia de pared con pseudoebullición en flujos supercríticos}. Mediante
simulación numérica directa con termodinámica de gas real, caractericé los regímenes de
flujo de fluidos transcríticos a alta presión en microconductos [J4]. También contribuí a una
base de datos abierta de flujos en canal microconfinado en régimen laminar, de transición y
turbulento [J3, D2], y colaboré en el análisis de los experimentos correspondientes con
métodos basados en datos [J7]. Para que estas simulaciones fueran fiables y asequibles,
establecí criterios de resolución para la DNS de turbulencia de pared transcrítica [J5] y
estudié unidades mínimas de flujo [J6]. Recientemente expliqué cómo la flotabilidad genera
flujos secundarios y redistribuye la cantidad de movimiento en la turbulencia transcrítica
en conductos [J2] (\emph{Editor's Pick} de \emph{Physics of Fluids}). Este trabajo constituye
la base de la \textbf{línea A} (turbulencia multifísica, compresible y supercrítica).

Todas estas simulaciones se apoyan en mi trabajo en la \textbf{línea C}. He liderado el
desarrollo de RHEA-X, la versión distribuida y residente en GPU del solver de código
abierto RHEA (MPI + OpenACC/CUDA) [J1]. Lo optimicé en los MareNostrum Hackathons del BSC
y lo he utilizado en asignaciones competitivas en MareNostrum~5, Finisterrae~III y
sistemas nacionales canadienses (más de 10 millones de horas de núcleo/GPU en total).

En 2026 inicié la \textbf{línea B} en el CTR Summer Program de la Universidad de Stanford,
con el proyecto \emph{inferencia causa--efecto en eventos extremos de flujo de calor en la
pared en flujos supercríticos}. El estudio lagrangiano de eventos termofluidodinámicos
extremos (RES IM-2026-1-0034) y el conjunto de datos para aprendizaje automático BLASTNet
[D1] lo amplían. Continuaré este trabajo en la University of British Columbia entre junio y
agosto de 2027. Este plan se construye directamente sobre estas tres bases. (Las
referencias [J*], [C*], [D*] siguen la numeración del HAP.)

\subsection{Motivación y objetivos de investigación}
Los fluidos supercríticos y transcríticos son clave en varias tecnologías de la transición
energética: ciclos de potencia de CO$_2$ supercrítico, refrigeración de motores cohete de
propulsante líquido y de estatorreactores (\emph{scramjets}), almacenamiento e inyección de
hidrógeno a alta presión, y reactores refrigerados por agua supercrítica. Cerca de la línea
pseudocrítica, las propiedades termofísicas varían en órdenes de magnitud en intervalos de
temperatura muy pequeños. La turbulencia resultante está fuertemente acoplada con la
termodinámica y la transferencia de calor. Los eventos raros e intensos de esta turbulencia,
como el deterioro local de la transferencia de calor, los picos de presión o la inversión del
flujo, determinan los márgenes de seguridad y los límites de diseño. Sin embargo, sus causas
se conocen poco y son difíciles de predecir. Las simulaciones de alta fidelidad pueden
resolver estos flujos, pero solo en grandes supercomputadores basados en GPU, y los datos
que generan requieren nuevas herramientas de análisis.

Este plan aborda el problema a lo largo de las tres líneas de la convocatoria:

\begin{description}[leftmargin=0.6cm]
\item[Línea A -- Flujos turbulentos multifísicos; fluidos compresibles y supercríticos.]
  \begin{itemize}
  \item[O1.] Generar bases de datos de alta fidelidad (DNS/LES con resolución de pared) de
  turbulencia de pared y de cizalladura libre a presión supercrítica, con termodinámica y
  propiedades de transporte de gas real. Incluirán canales y tuberías calentados y enfriados,
  así como capas de mezcla y chorros transcríticos.
  \item[O2.] Caracterizar el acoplamiento multifísico entre turbulencia, transferencia de
  calor y no idealidades termodinámicas (pseudoebullición, compresibilidad, flotabilidad), y
  su efecto en el deterioro y la mejora de la transferencia de calor.
  \end{itemize}
\item[Línea B -- Predicción basada en datos de flujos de ingeniería; eventos extremos; causalidad.]
  \begin{itemize}
  \item[O3.] Identificar los mecanismos causales que originan los eventos extremos (p.~ej.,
  picos de flujo de calor y de esfuerzo cortante en la pared) mediante causalidad basada en
  teoría de la información y estadística condicionada. Partirá del trabajo iniciado en el
  CTR-Stanford en 2026 y continuará en la UBC en 2027.
  \item[O4.] Desarrollar modelos basados en datos e informados por la causalidad que
  predigan estos eventos extremos y otras magnitudes de interés en ingeniería: precursores,
  modelos de orden reducido y modelos sustitutos basados en redes neuronales.
  \end{itemize}
\item[Línea C -- Modelización y computación; programación distribuida (CPU) y acelerada (GPU).]
  \begin{itemize}
  \item[O5.] Desarrollar y mantener solvers de código abierto y rendimiento portable para
  flujos compresibles de gas real en supercomputadores híbridos CPU--GPU (MPI +
  OpenACC/CUDA), con análisis de datos \emph{in situ} y acoplamiento con aprendizaje
  automático. Se basará en RHEA-X [J1], cuyo desarrollo lidero.
  \item[O6.] Desarrollar modelos de submalla y de pared para turbulencia supercrítica,
  coherentes con la termodinámica de gas real, entrenados y validados con las bases de datos
  de O1.
  \end{itemize}
\end{description}
\TODO{Adaptar los objetivos a tu experiencia real y mantener como máximo 2 por línea.}

\subsection{Paquetes de trabajo}
\textbf{PdT0 -- Recursos, creación de grupo y gestión (A1--A6).}
Solicitud de financiación competitiva y de tiempo de cálculo
(apartado~\ref{sec:resources}). Captación y dirección de estudiantes de doctorado y de
máster. Autoevaluación del tercer año y replanificación.

\textbf{PdT1 -- Turbulencia multifísica supercrítica (O1, O2; A1--A5).}
Configuración de casos DNS/LES con ecuaciones de estado de gas real (Peng--Robinson) y
modelos de transporte a alta presión. Validación con la literatura y con datos
experimentales, y publicación de bases de datos abiertas. \emph{Resultados esperados:} 2--3
artículos Q1 (p.~ej., J.\ Fluid Mech., Phys.\ Rev.\ Fluids, Int.\ J.\ Heat Mass Transfer)
y una base de datos abierta.

\textbf{PdT2 -- Causalidad y eventos extremos (O3; A1--A4).}
Aplicación de inferencia causal basada en teoría de la información a los datos del PdT1
para jerarquizar los mecanismos que originan los eventos extremos de transferencia de calor y
de cizalladura. Este paquete es el núcleo de la estancia en la UBC en 2027.
\emph{Resultados esperados:} 2 artículos Q1 y 1 tesis doctoral.

\textbf{PdT3 -- Predicción basada en datos (O4; A2--A6).}
Desarrollo de predictores y modelos sustitutos, informados por la causalidad, de eventos
extremos y magnitudes de ingeniería, con cuantificación de incertidumbre.
\emph{Resultados esperados:} 2 artículos Q1, modelos de código abierto y 1 tesis doctoral.

\textbf{PdT4 -- Solvers y modelos HPC (O5, O6; A1--A6).}
Portabilidad, optimización y escalado de los solvers en sistemas GPU europeos
preexaescala, y acoplamiento con entornos de aprendizaje automático. Desarrollo y
validación de modelos de submalla y de pared para gas real.
\emph{Resultados esperados:} versiones de software, 1--2 artículos (p.~ej., J.\ Comput.\
Phys., Comput.\ Phys.\ Commun.) y asignaciones de cálculo.

\textbf{PdT5 -- Movilidad, colaboraciones y difusión (A1--A6).}
Estancias de investigación, colaboraciones internacionales, congresos y divulgación
(apartados~\ref{sec:mobility}--\ref{sec:dissemination}).

\begin{figure}[H]
\centering
\begin{ganttchart}[
    hgrid, vgrid={*1{dotted}},
    x unit=0.55cm, y unit chart=0.55cm,
    time slot format=simple,
    title height=1,
    bar/.append style={fill=MSBlue!70},
    milestone/.append style={fill=red},
    bar label font=\footnotesize,
    group label font=\footnotesize\bfseries
  ]{1}{24}
  \gantttitle{A1}{4}\gantttitle{A2}{4}\gantttitle{A3}{4}\gantttitle{A4}{4}\gantttitle{A5}{4}\gantttitle{A6}{4}\\
  \ganttbar{PdT0 Recursos y grupo}{1}{24}\\
  \ganttbar{PdT1 Turbulencia supercrítica}{1}{20}\\
  \ganttbar{PdT2 Causalidad y extremos}{1}{16}\\
  \ganttbar{PdT3 Predicción basada en datos}{5}{24}\\
  \ganttbar{PdT4 Solvers y modelos HPC}{1}{24}\\
  \ganttbar{PdT5 Movilidad (UBC)}{2}{3}\ganttbar[inline]{}{6}{6}\ganttbar[inline]{}{10}{10}\\
  \ganttmilestone{Evaluación del 3.\textsuperscript{er} año}{12}\\
  \ganttmilestone{Defensa tesis 1}{18}\\
  \ganttmilestone{Defensa tesis 2}{23}
\end{ganttchart}
\caption{Cronograma por trimestres a lo largo de los seis años de contrato (A = año).
\TODO{alinear A1 con el inicio real del contrato (primer cuatrimestre posible, base 8.3) y
con las fechas reales de las estancias}}
\end{figure}

\subsection{Movilidad investigadora}\label{sec:mobility}
\begin{itemize}
\item \textbf{2026 (realizada)}: CTR Summer Program, Universidad de Stanford (1 mes, con
certificado), sobre fluidos supercríticos, eventos extremos y causalidad. Proyecto:
\emph{Inferencia causa--efecto en eventos extremos de flujo de calor en la pared en flujos
supercríticos}. \TODO{colaboradores}
\item \textbf{Junio--agosto de 2027 (financiada y ya organizada)}: University of British
Columbia (Canadá). Investigación en turbulencia multifísica y causalidad. Anfitri\'on:
\TODO{p.~ej., Prof.~Mona Rahmani -- confirmar}; financiación: \TODO{programa}
\item \textbf{2027--2030 (prevista, años 1--3)}: estancias breves en grupos europeos que
complementan el plan. Por ejemplo, el grupo del Prof.\ Mario Di Renzo (Università del
Salento, Italia), sobre simulación en GPU de flujos compresibles, de alta entalpía y
multifísicos. \TODO{verificar la afiliación y añadir 1--2 grupos europeos más y sus temas}
\end{itemize}
Estas estancias refuerzan la internacionalización y darán lugar a publicaciones y
solicitudes conjuntas (p.~ej., MSCA Doctoral Networks, convocatorias EuroHPC).

\subsection{Recursos y estrategia de financiación}\label{sec:resources}
\begin{itemize}
\item \textbf{Nacional}: Plan Estatal -- Proyectos de Generación de Conocimiento (PID),
primero como miembro del equipo y después como investigador principal. Convocatorias Ramón
y Cajal y Consolidación Investigadora cuando sea elegible. Contratos predoctorales FPI/FPU
vinculados a los proyectos.
\item \textbf{Autonómica}: AGAUR SGR (con el grupo del departamento), ayudas predoctorales
FI y Doctorados Industriales con empresas.
\item \textbf{Europea}: ERC Starting Grant (objetivo: años 3--5), MSCA Doctoral Networks y
Postdoctoral Fellowships con acogida en la UPC, Horizonte Europa Clúster 5 (Clima, Energía
y Movilidad).
\item \textbf{Cálculo}: asignaciones competitivas en la RES (MareNostrum~5), EuroHPC JU y
socios internacionales.
\item \textbf{Industria y transferencia}: contratos con empresas del sector energético
(ciclos de sCO$_2$, hidrógeno) y de propulsión aeroespacial, con el apoyo de los servicios
de transferencia de la UPC.
\end{itemize}
\TODO{enumerar los proyectos en los que ya participas y las primeras solicitudes previstas, con fechas}

\subsection{Dirección de tesis doctorales y creación de grupo}
Mi objetivo es codirigir dos tesis doctorales durante el contrato y dirigir una tercera en
solitario a partir del cuarto año. Temas propuestos:
\begin{enumerate}
\item Mecanismos causales y predicción de eventos extremos de transferencia de calor en
turbulencia de pared supercrítica (PdT1--PdT2--PdT3).
\item Modelización de submalla y de pared para gas real en LES acelerada en GPU de flujos
transcríticos (PdT1--PdT4).
\item Modelos de orden reducido basados en datos para componentes termofluidodinámicos de
ciclos de potencia de CO$_2$ supercrítico (PdT3, con la industria).
\end{enumerate}
Los trabajos de fin de grado y de máster (TFG/TFM) alimentarán estas líneas
(apartado~\ref{sec:theses}).

\subsection{Colaboraciones, transferencia y difusión}\label{sec:dissemination}
\begin{itemize}
\item \textbf{Colaboraciones}: Departamento de Mecánica de Fluidos de la UPC; CTR-Stanford;
UBC; grupos europeos visitados en el PdT5. \TODO{añadir colaboradores actuales}
\item \textbf{Revistas}: J.\ Fluid Mech., Phys.\ Rev.\ Fluids, J.\ Comput.\ Phys.,
Int.\ J.\ Heat Mass Transfer, Int.\ J.\ Heat Fluid Flow, J.\ Supercrit.\ Fluids.
\item \textbf{Congresos}: APS-DFD, ETMM, TSFP, ECCOMAS/WCCM, EUROMECH.
\item \textbf{Ciencia abierta}: códigos de código abierto (GitHub, DOI de Zenodo) y bases
de datos abiertas, siguiendo los principios FAIR y los planes de gestión de datos de los
proyectos.
\end{itemize}

% =====================================================================
\section{Proyecto docente}

\subsection{Experiencia docente}
En la EEBE he impartido la teoría completa de \emph{Termodinámica y Transferencia de Calor}
(TTC) \TODO{grados, cursos, horas} y las sesiones de laboratorio de \emph{Mecánica de
Fluidos} (MF) \TODO{grados, cursos, horas}. En el cuatrimestre de primavera de 2027
impartiré la teoría de \emph{Ingeniería de Fluidos}. Estas asignaturas cubren el núcleo del
perfil docente de la convocatoria: termodinámica, transferencia de calor y mecánica de
fluidos, tanto en el nivel fundamental como en el aplicado. En el ICE de la UPC he
completado el curso \emph{M2.4 -- Atenea: curso de nivel medio} (2026), que aplico al diseño
de las asignaturas en el campus virtual (actividades condicionales, grupos de laboratorio,
control de asistencia). He realizado además varios cursos de
catalán (nivel actual B1) \TODO{confirmar}.

\subsection{Contexto: los grados de la EEBE}
La EEBE ofrece grados de 240 ECTS que comparten un bloque común inicial. En ese bloque,
\emph{Mecánica de Fluidos} (3.\textsuperscript{er} cuatrimestre) y \emph{Termodinámica y
Transferencia de Calor} (4.\textsuperscript{o} cuatrimestre) son obligatorias en todos los
grados de ingeniería \TODO{comprobar los cuatrimestres en los planes de estudios vigentes}.
Las asignaturas específicas de cursos posteriores se apoyan en ellas. El Máster
Universitario en Tecnologías Mecánicas incluye \emph{Mecánica de Fluidos Computacional}. La
docencia sigue el marco del EEES: carga de trabajo en ECTS (25--30 h/ECTS), aprendizaje
basado en competencias y evaluación continua.

\subsection{Asignaturas de la convocatoria y adecuación}
\begin{center}
\small
\begin{tabularx}{\textwidth}{@{}p{3.9cm}p{3.6cm}X@{}}
\toprule
\textbf{Asignatura} & \textbf{Titulación(es)} & \textbf{Mi contribución} \\
\midrule
Termodinámica y Transferencia de Calor & Grados en Ingeniería de la Energía, Mecánica, Biomédica, Electrónica Industrial y Automática, y Eléctrica & Teoría impartida (apartado~\ref{sec:ttc}); herramientas computacionales de propiedades; laboratorios de transferencia de calor \\
Termodinámica & Grado en Ingeniería de Materiales & Mismos fundamentos que TTC; propiedades de fluidos reales \\
Mecánica de Fluidos & Todos los grados de la EEBE & Laboratorio impartido; teoría y problemas \\
Ingeniería de Fluidos & Grado en Ingeniería Mecánica & Teoría a partir de la primavera de 2027 (apartado~\ref{sec:ef}) \\
Ingeniería Térmica & Grado en Ingeniería Mecánica & Intercambiadores de calor y sistemas térmicos; vínculo con la transferencia de calor supercrítica \\
Centrales Eléctricas y Energías Renovables & Grado en Ingeniería Eléctrica & Ciclos Rankine y Brayton, ciclos de sCO$_2$, energías renovables \\
Generación Termofluidodinámica & Grado en Ingeniería de la Energía & Turbomáquinas y generación de energía térmica \\
Mecánica de Fluidos Computacional & Máster Universitario en Tecnologías Mecánicas & CFD, HPC (MPI/CUDA) y análisis de flujos (apartado~\ref{sec:cfm}) \\
\bottomrule
\end{tabularx}
\end{center}

\subsection{Asignatura desarrollada: Termodinámica y Transferencia de Calor (TTC)}\label{sec:ttc}
\NOTE{Adaptado del apartado 4.5 del proyecto de Profesor Agregado; reescribir con tus propias palabras y contrastar con la guía docente vigente de TTC.}

\textbf{Requisitos previos.} Haber superado Física II y Cálculo Numérico--Ecuaciones
Diferenciales; cursar Mecánica de Fluidos en paralelo o con anterioridad.

\textbf{Objetivos de aprendizaje.} Al finalizar la asignatura, el estudiantado debe ser capaz de:
\begin{itemize}
\item enunciar y aplicar los principios cero, primero y segundo de la termodinámica a sistemas cerrados y abiertos;
\item modelizar componentes (cilindro--pistón, turbinas, compresores, intercambiadores de
calor), representar procesos en diagramas termodinámicos y evaluar propiedades de sustancias reales;
\item analizar ciclos de potencia y de refrigeración (Rankine, Brayton, compresión de vapor) y calcular su rendimiento;
\item distinguir la conducción, la convección y la radiación, y aplicarlas a problemas de aislamiento, refrigeración e intercambio de calor.
\end{itemize}
Competencia transversal: aprendizaje autónomo (nivel 2).

\textbf{Temario.}
(1) Fundamentos: sistemas, temperatura, gas ideal, propiedades de sustancias puras.
(2) Primer principio: trabajo, energía interna, calor, entalpía, calores específicos, procesos politrópicos, sistemas cerrados y abiertos.
(3) Segundo principio: entropía, irreversibilidades, máquinas térmicas, ciclo de Carnot, rendimientos isentrópicos, ciclos Brayton, Rankine y de refrigeración.
(4) Conducción: ecuación del calor, paredes planas y cilíndricas, resistencias térmicas, coeficiente global.
(5) Convección: natural y forzada, flujos internos y externos, correlaciones empíricas.
(6) Radiación: cuerpo negro y cuerpo gris, ley de Kirchhoff, factores de visión, intercambio radiativo.

\textbf{Metodología.} Se combinan clases expositivas participativas con el \emph{aula
invertida}. Cada semana, una sesión de una hora introduce el bloque siguiente y el material
que debe estudiarse (vídeos breves, apuntes, ejemplos resueltos). La sesión de dos horas se
dedica a la resolución de problemas en grupos reducidos. El estudiantado utiliza
herramientas computacionales de propiedades (p.~ej., la librería de código abierto CoolProp
en cuadernos de Python) para explorar el comportamiento de fluidos reales y de ciclos. Un
miniproyecto vincula la asignatura con la tecnología actual, p.~ej., un ciclo de CO$_2$
supercrítico.

\textbf{Dedicación (6 ECTS, 150 h).}
\begin{center}\small
\begin{tabular}{@{}llrr@{}}
\toprule
& Actividad & Horas & \% \\ \midrule
Dirigida & Teoría y ejemplos & 30 & 20 \\
         & Resolución de problemas en grupo & 15 & 10 \\
         & Laboratorio & 15 & 10 \\
Autónoma & Estudio, problemas y proyecto & 90 & 60 \\ \bottomrule
\end{tabular}
\end{center}

\textbf{Evaluación.} Cuestionarios previos en Atenea (15\,\%); entregables en el aula
(15\,\%); informes de laboratorio (20\,\%); examen parcial (20\,\%); examen final (30\,\%).
\TODO{contrastar con la normativa de evaluación vigente de TTC en la EEBE}

\textbf{Prácticas de laboratorio.} Termometría, calorimetría, conducción, procesos y ciclos, radiación.

\textbf{Bibliografía.} Y.\ A.\ \c{C}engel, M.\ A.\ Boles, \emph{Termodinámica},
McGraw-Hill. Y.\ A.\ \c{C}engel, A.\ J.\ Ghajar, \emph{Transferencia de calor y masa:
fundamentos y aplicaciones}, McGraw-Hill. M.\ J.\ Moran, H.\ N.\ Shapiro, \emph{Fundamentos
de termodinámica técnica}, Reverté.

\textbf{Planificación (15 semanas).} Semanas 1--2: tema 1; semanas 3--6: tema 2; semanas
7--9: tema 3 (examen parcial en la semana 8); semanas 10--11: tema 4; semanas 12--13: tema
5; semanas 14--15: tema 6. Sesiones de laboratorio cada dos semanas.

\subsection{Ingeniería de Fluidos}\label{sec:ef}
Impartiré la teoría de esta asignatura a partir de la primavera de 2027.
\TODO{Copiar el temario de la guía docente vigente (p.~ej., flujo en conductos y redes,
bombas y turbomáquinas, semejanza, cavitación, sistemas de fluidos). Describir la
metodología prevista: casos de estudio de instalaciones reales, proyectos de selección de
bombas y solvers sencillos de redes en Python. Describir la evaluación.}

\subsection{Asignatura de máster: Mecánica de Fluidos Computacional}\label{sec:cfm}
La asignatura (aproximadamente 1,5~h de teoría + 1,5~h de laboratorio por semana) abarca
Python para computación científica, métodos numéricos (EDO, EDP, álgebra lineal), un solver
de Navier--Stokes incompresible/compresible, computación de altas prestaciones (OpenMP, MPI,
CUDA) y análisis computacional de flujos (estabilidad, POD, DMD). La evaluación combina
entregables, un proyecto de curso por parejas y un examen final. Mi experiencia
investigadora se relaciona directamente con las partes de HPC y de flujo compresible.
Propongo:
\begin{itemize}
\item una sesión de laboratorio en GPU que acelere el ejemplo de Jacobi/Poisson de la
asignatura con CUDA u OpenACC y mida la aceleración y la escalabilidad;
\item ampliar el tema de análisis de flujos con análisis de series temporales basado en
datos y en causalidad (p.~ej., el caso de la cavidad con tapa móvil), vinculado a la línea de
investigación B;
\item proyectos de curso que puedan convertirse en TFM dentro del plan de investigación.
\end{itemize}
\TODO{acordar la contribución concreta con los coordinadores de la asignatura}

\subsection{Innovación docente}
\begin{itemize}
\item Aula invertida y vídeos breves compartidos entre el profesorado de la asignatura, para
homogeneizar los contenidos entre grupos.
\item Cuadernos computacionales (Python, CoolProp, CFD sencilla) en TTC, MF e Ingeniería de
Fluidos, para acercar las herramientas de fluidos reales y de simulación al nivel de grado.
\item Aprendizaje basado en proyectos sobre casos de la transición energética (ciclos de
sCO$_2$, hidrógeno, bombas de calor).
\item Uso responsable de la inteligencia artificial generativa en la resolución de problemas
y en los informes, de acuerdo con las directrices de la UPC.
\item Formación continua en el ICE, continuando el itinerario iniciado con el curso de Atenea (2026), con prioridad para metodologías activas y evaluación, y participación en proyectos de
innovación docente.
\end{itemize}

\subsection{Trabajos de fin de grado, de máster y tesis doctorales}\label{sec:theses}
\textbf{TFG}: análisis termodinámico de un ciclo Brayton de CO$_2$ supercrítico; estudio
experimental y numérico de un intercambiador de calor; aceleración en GPU de un solver de la
ecuación del calor.\\
\textbf{TFM}: LES de flujo calentado en tubería a presión supercrítica; análisis causal de
eventos de transferencia de calor en bases de datos de turbulencia; modelo sustituto basado
en redes neuronales del flujo de calor en la pared.\\
\textbf{Tesis doctorales}: temas indicados en el proyecto de investigación.

\subsection{Compromiso lingüístico}
Impartiré docencia en catalán, castellano e inglés, de acuerdo con el plan lingüístico de la
EEBE. Actualmente tengo el nivel B1 de catalán \TODO{confirmar} y he realizado varios
cursos de catalán \TODO{enumerar}. He solicitado la exención temporal prevista en el
art.~4.2 de la Normativa de acreditación del conocimiento de catalán del PDI de la UPC (Acuerdo
CG/2025/06/29) y, al incorporarme, me inscribiré en el itinerario de formación en catalán del
PDI del Servicio de Lenguas y Terminología (SLT). Mi objetivo es obtener el certificado C1
(itinerario del SLT o examen CIFALC) en los dos primeros años, muy por delante de la
evaluación del tercer año y del plan de estabilización, que lo exigen.

\subsection{Igualdad e inclusión}
Aplicaré el Plan de Igualdad de Género de la UPC en la docencia y en la investigación. Esto
implica utilizar un lenguaje y unos ejemplos inclusivos, buscar el equilibrio en la
captación de estudiantes e investigadores, atender al estudiantado con necesidades
específicas y apoyar a las mujeres en la ingeniería mediante actividades de divulgación y
mentoría.

% =====================================================================
\section{Hitos del tercer año (autoevaluación)}
\begin{itemize}
\item Docencia: docencia regular en al menos 3 asignaturas del perfil; material de aula
invertida de TTC disponible; Ingeniería de Fluidos consolidada; formación en el ICE
completada; nivel C1 de catalán.
\item Investigación: 4--6 artículos Q1 en las tres líneas; al menos una solicitud de
proyecto competitivo como investigador principal o coinvestigador principal; al menos una
asignación de cálculo; 1--2 doctorandos incorporados.
\item Internacionalización: estancia en la UBC (2027) y al menos una estancia europea
realizadas, con publicaciones conjuntas.
\end{itemize}

\end{document}
